\documentclass[11pt,a4paper]{article}
\usepackage[margin=2.4cm]{geometry}
\usepackage{booktabs}
\usepackage{parskip}

\title{The ten-name ranking on verified fills,\\
\large and what goes into the daily book}
\author{Bayesian Capital}
\date{9 August 2026}

\begin{document}
\maketitle

\section*{Basis}

Every name on the identical instrument: residual OU sigma, verified same-day fills
from 5-minute bars, single split at 2025-05-23 (no multi-fold slicing --- each tested
number stands on 30--120 verified fills). ``Ref'' is the full-sample-provenance vector
(deployed for incumbents; the diversifier-study fits for GM/VLO/CF; the NVDA daily
vector; AVGO fitted here on the full sample, flagged) scored on the tested half ---
all ten carry the same lookahead flavour, so they are comparable with each other.
A and B are the train-half-only fits, the genuinely out-of-sample numbers.
Candidates' daily bars are derived from the 5-minute regular-hours data; GM, VLO and
CF reproduce the diversifier study's full-sample figures exactly (50.3\%, 64.2\%,
46.9\%), confirming the basis.

\section*{The ranking (tested half, annualised, verified)}

\begin{center}
\begin{tabular}{llrrrcrrrr}
\toprule
 & & ref TEST & A TEST & B TEST & planned range & plan & buys/yr & maxDD & AI $\beta$ \\
\midrule
RKLB & inc & 159.7\% & 12.1\% & 34.3\% & 156--160\% & 161\% & 119 & 40.0\% & 0.79 \\
VRT  & inc & 157.4\% & 63.3\% & 50.0\% & $-5$--157\% & 58\% &  95 & 34.0\% & 1.19 \\
MU   & inc & 148.5\% & 113.4\% & 73.7\% & 7--149\% & 65\% & 84 & 31.7\% & 1.14 \\
VLO  & $+$ & 137.9\% & 41.2\% & 30.2\% & 12--138\% & 64\% &  49 & 10.7\% & 0.11 \\
TSM  & inc &  80.7\% & 58.4\% & 76.4\% & 33--81\% & 56\% &  99 & 16.7\% & 0.69 \\
GM   & $+$ &  56.5\% & 42.7\% & 31.8\% & 45--57\% & 50\% &  56 & 11.3\% & 0.17 \\
CF   & $+$ &  54.7\% & 52.9\% & 42.4\% & 39--55\% & 47\% &  31 & 17.4\% & $-0.01$ \\
AVGO & $+$ &  53.7\% & 21.7\% & 35.2\% & 49--54\% & 51\% &  60 & 22.5\% & 0.74 \\
NVDA & $+$ &  46.9\% & 49.0\% & 28.3\% & 33--47\% & 40\% &  38 & 12.4\% & 0.70 \\
VST  & inc &  46.2\% & 25.4\% & 32.6\% & 46--53\% & 50\% &  80 & 16.3\% & 0.98 \\
\bottomrule
\end{tabular}
\end{center}

\textbf{Planned range} is the reference vector's verified annualised return on the two
disjoint half-samples --- the train half contains the Jan--Apr 2025 AI crash, the test
half is the bull run --- so the range brackets what a regime can do to the name;
\textbf{plan} is the full-sample through-cycle figure. Both carry the vector's
full-sample fit provenance. The A/B columns are the independent fragility check: where
they collapse far below the range (RKLB 12\%, AVGO 22\%), the plan leans entirely on
one specific parameter vector that cannot be recovered from half the data, and the low
end should be planned on. A wide range (VRT, MU, VLO) means the name's earnings are
regime-driven --- the honest expectation for any single year is anywhere inside it.
Book correlations (tested half): RKLB 0.94, MU 0.59, VRT 0.48, TSM 0.47, NVDA 0.43,
VST 0.35, AVGO 0.31, GM 0.24, VLO 0.04, CF $-0.01$.

Book effect of the additions (held construction, common calendar):

\begin{center}
\begin{tabular}{lrrrr}
\toprule
 & full \%/yr & full maxDD & test \%/yr & test maxDD \\
\midrule
five (current)        & 84.8\% & 39.3\% & 122.4\% & 29.6\% \\
eight ($+$GM/VLO/CF)  & 71.4\% & 29.7\% & 101.9\% & 20.1\% \\
\bottomrule
\end{tabular}
\end{center}

Jan--Apr 2025 stress (the one AI drawdown in sample): five $-18.7\%$, eight
$-11.8\%$ on verified strategy equity.

\section*{Decision}

\textbf{Add GM, VLO and CF.} Tested 55--138\% with AI beta 0.0--0.17, book
correlation 0.0--0.24 and single-name drawdowns of 11--17\% against the incumbents'
17--40\%. The additions cost about 13pp of book return and buy 10pp of drawdown and
a third off the stress loss --- the same trade the June study measured, now confirmed
on the verified single-split instrument. CF deserves note as the most robust name in
the entire table: its train-only fits (52.9\%, 42.4\%) sit within a few points of its
full-sample reference (54.7\%) --- the only name where fresh fitting transports ---
and it is the one true zero-beta name.

\textbf{Do not add NVDA or AVGO to the daily book.} (The weekly book has been
withdrawn, so this is judged purely on the numbers.) Under the daily model they earn
47--54\% tested and plan at 40--51\% --- below or level with every incumbent except
VST --- while carrying AI betas of 0.70--0.74 and book correlations of 0.3--0.4. At
equal weight a new name takes capital from what it dilutes: these two would lower the
book's return without lowering its risk, which is the opposite of what an addition is
for. AVGO's fragility check is also the second-worst in the table (A tested 21.7\%
against its 49--54\% range), and its reference vector was fitted in this session.

\textbf{VST is the weakest slot in the book.} Lowest tested return (46.2\%), highest
AI beta (0.98), and its fresh fits collapse (25--33\%). Nothing here requires
removing it, but if the eight-name allocation ever needs a funding source or the
book needs to shed a name, VST is the candidate --- flagged, not actioned.

\textbf{Planning figures} stay at the June levels --- GM 33\%, VLO 35\%, CF 41\% ---
which the new tested numbers are consistent with. VLO's 137.9\% tested half is one
energy rally (its train half earned 11.5\%) and must not be planned on.

Weights on admission: eight names at 12.5\% each, floor and cap both 0.125, per the
open implementation item in the handover.

Reproduction: \texttt{premarket\_study/fresh\_opt\_cands.py} (candidate sweep),
\texttt{rank\_book.py} (this table); derived data untracked.

\end{document}
